\documentclass[11pt,a4paper]{article}
\usepackage[spanish]{babel}
\usepackage[utf8]{inputenc}
\usepackage[T1]{fontenc}
\usepackage{lmodern,amsmath,amssymb,geometry,enumitem}
\geometry{margin=2.5cm}
\setlength{\parindent}{0pt}
\setlength{\parskip}{5pt}
\title{Entregables Estructuras Algebraicas}
\author{}
\date{}
\newcommand{\Z}{\mathbb Z}
\newcommand{\F}{\mathbb F}
\newcommand{\C}{\mathbb C}
\newcommand{\ej}[1]{\par\medskip\textbf{Ejercicio #1.}\enspace}
\begin{document}
\maketitle
\section{Operaciones, axiomas y órdenes}
\ej{1} En el intervalo $G=(-1,1)$ se define la operación
\[x*y=\frac{x+y}{1+xy}.\]
¿Es $(G,*)$ un grupo?

\newpage

\ej{2} Inventa un grupo infinito en el que cada elemento tenga orden finito (por ejemplo, 2).

\vspace{5cm}

\ej{3}
\begin{enumerate}[label=(\alph*)]
\item Demuestra que un grupo $G$ es abeliano si y sólo si $(ab)^2=a^2b^2$ para todos $a,b\in G$.
\item Demuestra que $G$ es abeliano si y sólo si $f(x)=x^{-1}$ es un automorfismo de $G$.
\end{enumerate}

\newpage

\ej{4} Demuestra que si un grupo tiene orden par, alguno de sus elementos, además del neutro, es su propio inverso.

\vspace{5cm}

\ej{6} Describe el grupo de isometrías de un rectángulo no cuadrado.

\vspace{5cm}

\ej{12} ¿Qué subgrupo de $\C^*$ tiene 8 elementos? ¿Cuáles son sus generadores?

\vspace{5cm}

\ej{13} ¿A qué grupo es isomorfo $\Z_9^*$?

\newpage

\ej{14} Demuestra que en el juego del 15 no se pueden intercambiar la 14 y la 15 sin alterar las demás fichas.
\[\begin{array}{cccc}1&2&3&4\\5&6&7&8\\9&10&11&12\\13&14&15&\phantom{0}\end{array}\]

\vspace{5cm}

\ej{16} ¿Existe alguna combinación de rotaciones del cubo de Rubik que, aplicada cierto número de veces, pueda generar todas las posiciones?

\vspace{5cm}

\ej{18} Encuentra todos los subgrupos normales de $S_4$.
\end{document}
